\chapter[Comprehensive User Guide: Instructions for Using the Template]{Comprehensive User Guide Instructions for Using the Template}
\label{cp:user-guide}

If you plan to use this template, please read this chapter carefully. It provides all the information you need to effectively use the template, including the mandatory modifications (\textit{e.g.}, title, subtitle, author information) and other configurations that, while not highly recommended, are optional. The template comprises various directories and files, including a total of seven distinct directories and dozens of files. Among these, the most important are \texttt{EPFLMain.tex} and \texttt{EPFLThesis.cls}. Below, \autoref{tab:file-structure} presents the different directories available, along with their descriptions and a check-mark indicating whether you need to access the directory to make changes. Of course, the check mark indicates that you can make changes to the content, while a hyphen signifies that you should not modify it.

\begin{table}[!htpb]
    \setlength{\extrarowheight}{2pt}
    \caption[Directory structure and file organisation]{Overview of the directory structure in this template.}
    \label{tab:file-structure}
    \begin{tabularx}{\textwidth}{lcX}
        \toprule
        \\[-1.5\normalbaselineskip]
        \textbf{Directory} & \textbf{Modifiable} & \textbf{Description} \\ [0em]
        \midrule
        \textit{Bibliography} & $\checkmark$ & This folder contains the bibliography file used to manage references throughout the document. \\
        \textit{Chapters} & $\checkmark$ & Individual chapters of the thesis are organised in this directory, making it easy to work on sections separately. \\
        \textit{Code} & $\checkmark$ & Code examples and relevant scripts are stored here, supporting the content of the thesis. \\
        \textit{Configurations} & - & All configuration files required for the template, such as layout and style settings, are placed in this directory. \\
        \textit{Figures} & $\checkmark$ & All figures and images referenced within the document are stored in this folder for easy access and management. \\
        \textit{Matter} & - & The front matter of the document, including the cover page, copyright statement, and glossary, is assembled in this directory. \\
        \textit{Metadata} & $\checkmark$ & This folder holds the metadata file, where key document details such as the author, title, and supervisor can be customised. \\
        \bottomrule
    \end{tabularx}
\end{table}

It is crucial to note that the files are organised according to a specific naming convention, which must be \textbf{respected} and \textbf{maintained}. The naming convention consists of an ascending two-digit numeric value, followed by a hyphen, and then the file name in capital letters. The name should always aim to be a single word. If more than one word is necessary, they should be separated by a hyphen and capitalised.

\begin{block}[note]
While \autoref{tab:file-structure} indicates that the \textit{Matter} directory is not modifiable, two files within that directory should be altered when necessary: \texttt{05-Glossary.tex} and \texttt{06-Acronyms.tex}. Although the names are fairly self-explanatory, these files should contain the glossary and acronyms entries, respectively.
\end{block}

The two files mentioned earlier, \texttt{EPFLMain.tex} and \texttt{EPFLThesis.cls}, should be used with caution. The main file, as the name suggests, is the master file where you will add the necessary chapters to be included in your work. The class file, on the other hand, requires even more caution, and it is not recommended to alter it.

\section{Template and Class Options}
\label{sec:class-options}
The first thing you need to do is specify the options within the \texttt{EPFLMain.tex} file. How do you do that? It's simple. At the top of the file, you will find a \texttt{documentclass} command that loads the custom class for this template. In this call, you can pass the options you need. The available options, presented in a key-argument style, are listed in \autoref{tab:template-options}.

{
\setlength{\extrarowheight}{-1.75pt}
\begin{xltabular}{\textwidth}{lX}
\caption{Class options supported by the template.}
\label{tab:template-options} \\
%
\toprule 
\multicolumn{1}{l}{\textbf{Options}} & \multicolumn{1}{l}{\textbf{Description}} \\ 
\midrule
\endfirsthead
%
\multicolumn{2}{c}%
{{\textit{\bfseries Table \thetable\ continued from previous page.}}} \\
%
\toprule 
\multicolumn{1}{l}{\textbf{Options}} & \multicolumn{1}{l}{\textbf{Description}} \\ 
\midrule
\endhead
%
\bottomrule
\addlinespace[1mm]
\multicolumn{2}{r}%
{{\textit{Continued on the next page.}}} \\
\endfoot
\bottomrule
\endlastfoot

\textbf{language=OPT} & \textbf{Language preference selection.} \\
\footnotesize{\textit{portuguese, english, spanish}} & \footnotesize{\textit{$\Rightarrow$ Default: language=english}} \\[0.85em]
        
\textbf{chapterstyle=OPT} & \textbf{Selection of a chapter design style.} \\
\multirow[t]{2}{*}{\footnotesize{\textit{classic, modern, fancy}}} & \footnotesize{\textit{$\Rightarrow$ Default: chapterstyle=classic}} \\
& \footnotesize{\textit{This option modifies the appearance of the chapter, including its title and numbering style. Explore the available styles and apply the one you prefer.}} \\[1.70em]

\textbf{coverstyle=OPT} & \textbf{Choosing a style for the cover.} \\
\multirow[t]{3}{*}{\footnotesize{\textit{classic, bw}}} & \footnotesize{\textit{$\Rightarrow$ Default: coverstyle=classic}} \\
& \footnotesize{\textit{classic $\rightarrow$ EPFL cover with red gradient background.}} \\
& \footnotesize{\textit{bw $\rightarrow$ Black and white cover variant.}} \\

\textbf{docstage=OPT} & \textbf{Choosing a stage for you document.} \\
\multirow[t]{3}{*}{\footnotesize{\textit{final, working}}} & \footnotesize{\textit{$\Rightarrow$ Default: docstage=final}} \\
& \footnotesize{\textit{final $\rightarrow$ Assumes this is the final version of the document.}} \\
& \footnotesize{\textit{working $\rightarrow$ It assumes the document is a work in progress.}} \\[.3em]

\textbf{media=OPT} & \textbf{Project media type.} \\
\multirow[t]{3}{*}{\footnotesize{\textit{paper, screen}}} & \footnotesize{\textit{$\Rightarrow$ Default: media=paper}} \\
& \footnotesize{\textit{paper $\rightarrow$ Blank pages will appear between sections.}} \\
& \footnotesize{\textit{screen $\rightarrow$ Blank pages will not appear between sections.}} \\[.3em]

\textbf{linkcolor=OPT} & \textbf{Main theme color.} \\
\multirow[t]{2}{*}{\footnotesize{\textit{color}}} & \footnotesize{\textit{$\Rightarrow$ Default: linkcolor=red!45!black}} \\
& \footnotesize{\textit{This option requires a valid color name. Refer to the xcolor manual (subsection 4.2) to select a valid color.}} \\[.3em]

\textbf{bookprint=OPT} & \textbf{For book printing.} \\
\multirow[t]{2}{*}{\footnotesize{\textit{true, false}}} & \footnotesize{\textit{$\Rightarrow$ Default: bookprint=false}} \\
& \footnotesize{\textit{This option adds a binding margin on odd-numbered pages to allow for printing, as it increases the left margin.}} \\[.3em]

\textbf{aiacknowledgement=OPT} & \textbf{AI acknowledgement print.} \\
\multirow[t]{2}{*}{\footnotesize{\textit{true, false}}} & \footnotesize{\textit{$\Rightarrow$ Default: aiacknowledgement=true}} \\
& \footnotesize{\textit{This option adds a section intended for the user to insert their acknowledgement of AI usage.}} \\[.3em]

\textbf{listprefix=OPT} & \textbf{Add a prefix to LoF and LoT.} \\
\multirow[t]{2}{*}{\footnotesize{\textit{true, false}}} & \footnotesize{\textit{$\Rightarrow$ Default: listprefix=false}} \\
& \footnotesize{\textit{This options adds the prefix ``Figure'' or ``Table'' to both the LoF and LoT when enabled.}} \\
\end{xltabular}
}

After setting the desired options in the main class, you are all set to personalise the metadata. To learn how, simply refer to \autoref{sec:metadata}.

\section{Metadata Customisation}
\label{sec:metadata}
While options like language can be passed as arguments to the main class, other options, such as author and title, must be defined manually. Since this template supports a wide range of metadata options, a dedicated file is provided for this purpose. The file at \texttt{Metadata/Metadata.tex} lists metadata variables, with comments on whether they are mandatory. Comment out the variables to omit them. \autoref{tab:metadata} includes all metadata variables, their GET command, and if they are mandatory. The GET command automatically retrieves the information from the stored variable.

\begin{longtable}[c]{llc}
\caption{Metadata variables within the template.}
\label{tab:metadata} \\
\toprule
\textbf{Variable} & \textbf{Macro Commands} & \textbf{Mandatory} \\ \midrule
\endfirsthead
%
\multicolumn{3}{c}%
{{\textit{\bfseries Table \thetable\ continued from previous page.}}} \\
%
\toprule
\textbf{Variable} & \textbf{Macro Commands} & \textbf{Mandatory} \\ \midrule
\endhead
%
\bottomrule
%
\addlinespace[1mm]
\multicolumn{3}{r}%
{{\textit{Continued on the next page.}}} \\
\endfoot
%
\bottomrule
%
\endlastfoot
%
Title            & \verb|\GetTitle|         & $\checkmark$ \\
Subtitle         & \verb|\GetSubtitle|      & $\checkmark$ \\
University       & \verb|\GetUniversity|    & $\checkmark$ \\
School           & \verb|\GetSchool|        & $\checkmark$ \\
Department       & \verb|\GetDepartment|    & $\checkmark$ \\
Degree           & \verb|\GetDegree|        & $\checkmark$ \\
Course           & \verb|\GetCourse|        & -            \\
Local and date   & \verb|\GetDate|          & $\checkmark$  \\ 
Academic year    & \verb|\GetAcademicYear|  & $\checkmark$ \\ 

Thesis type (\scriptsize{\textit{Dissertation, Project or Internship}}) & \verb|\GetThesisType| & $\checkmark$ \\

First author name           & \verb|\GetFirstAuthor|        & $\checkmark$ \\
First author identification & \verb|\GetFirstAuthorNumber|  & $\checkmark$ \\ 

Second author name           & \verb|\GetSecondAuthor|          & - \\
Second author identification & \verb|\GetSecondAuthorNumber|    & - \\ 

Third author name           & \verb|\GetThirdAuthor|        & - \\
Third author identification & \verb|\GetThirdAuthorNumber|  & - \\ 

Supervisor name                  & \verb|\GetSupervisor|        & $\checkmark$ \\
Supervisor e-mail                & \verb|\GetSupervisorMail|    & $\checkmark$ \\
Supervisor title and affiliation & \verb|\GetSupervisorTitle|   & $\checkmark$ \\ 

Co-supervisor name                  & \verb|\GetCoSupervisor|       & - \\
Co-supervisor e-mail                & \verb|\GetCoSupervisorMail|   & - \\
Co-supervisor title and affiliation & \verb|\GetCoSupervisorTitle|  & - \\ 

Second co-supervisor name                   & \verb|\GetSecCoSupervisor|      & - \\
Second co-supervisor e-mail                 & \verb|\GetSecCoSupervisorMail|  & - \\
Second co-supervisor title and affiliation  & \verb|\GetSecCoSupervisorTitle| & - \\
\end{longtable}

If you want to add more metadata options, open an issue in the \href{https://github.com/pmdlt/epfl-thesis}{GitHub repository}.

\section{Custom Chapter Insertion}
As stated before, to use this template, you need to do three things: set the appropriate options in the document class (see \autoref{sec:class-options}), update the document metadata (see \autoref{sec:metadata}), and create and import your custom chapters. To create and import a custom chapter, follow these steps: \(i\) create a TeX file in the Chapters directory that follows the predefined naming convention and \(ii)\) include it in the main file using the command \verb|\include{CHAPTER}|. And voilà, your first chapter is ready!

\section{Custom Commands}
Within this template, some custom commands are also available for your use. For example, if you are writing your thesis and want to add a to-do note, you can easily insert a block with the option \verb|todo|, as follows: \verb|\begin{block}[todo]|. This will insert a to-do block with a style similar to Markdown. Other available options are: \verb|tip|, \verb|warning|, and \verb|note|. Below is a visual example for each one.

\vspace{.875em}
\begin{tcbraster}[
    raster columns=2, 
    raster equal height, 
    nobeforeafter, 
    raster column skip=1cm
]
\begin{block}[todo]
    This is a to-do block.
\end{block}
\begin{block}[tip]
    This is a tip block.
\end{block}
\end{tcbraster}

\begin{tcbraster}[
    raster columns=2, 
    raster equal height, 
    nobeforeafter, 
    raster column skip=1cm
]
\begin{block}[warning]
    This is a warning block.
\end{block}
\begin{block}[note]
    This is a note block.
\end{block}
\end{tcbraster}
\vspace{.875em}

\section{EPFL Brand Identity}
\label{sec:brand-identity}
This template follows the EPFL brand guidelines (\href{https://corp-id.epfl.ch}{corp-id.epfl.ch}) wherever a thesis document allows it. This section documents how the guidelines are applied and presents the reusable brand elements that ship with the template.

\subsection{Corporate Colors}
\label{sec:brand-colors}
The guidelines state that, when using colors, you must always choose the values of the official palette, and always balance colors with generous white space. The six corporate colors are predefined in the template (see \texttt{Configurations/01-Colors.sty}) and can be used anywhere with \verb|\textcolor{epflrouge}{...}| or any other \texttt{xcolor} command. \autoref{tab:epfl-colors} lists the palette with the official print and digital values.

\begin{table}[!htpb]
    \setlength{\extrarowheight}{2.5pt}
    \caption[EPFL corporate color palette]{The EPFL corporate color palette and its official values.}
    \label{tab:epfl-colors}
    \begin{tabularx}{\textwidth}{clllX}
        \toprule
        & \textbf{Name} & \textbf{Hex} & \textbf{Pantone / RAL} & \textbf{Macro} \\
        \midrule
        \epflswatch{epflrouge}     & Rouge     & \texttt{\#FF0000} & 485 C / 3020     & \texttt{epflrouge} \\
        \epflswatch{epflgroseille} & Groseille & \texttt{\#B51F1F} & 3546 CP / 3000   & \texttt{epflgroseille} \\
        \epflswatch{epfltaupe}     & Taupe     & \texttt{\#413D3A} & 405 CP / 8019    & \texttt{epfltaupe} \\
        \epflswatch{epflcanard}    & Canard    & \texttt{\#007480} & 7714 CP / 5025   & \texttt{epflcanard} \\
        \epflswatch{epflleman}     & Léman     & \texttt{\#00A79F} & 3272 CP / 5018   & \texttt{epflleman} \\
        \epflswatch{epflperle}     & Perle     & \texttt{\#CAC7C7} & 406 C / 7047     & \texttt{epflperle} \\
        \bottomrule
    \end{tabularx}
\end{table}

\subsection{Logo and Brand Architecture}
The EPFL logo always features on the top left side of documents, in Swiss red, primarily on clear backgrounds and images. The logo should remain red as much as possible, even on dark backgrounds; the white version on a red background is reserved for secondary uses. This template honours these rules on the cover, front page, and back page.

The openings of the letters ``E'' and ``F'' in the logo come from the proportions of the ends of the arm of the Swiss cross: one sixth wider than high. The same construction is available as the \verb|\epflsquare| element (\epflsquare), which the guidelines pair with organizational unit references. Unit references must always be separated from the EPFL logo, accompanied by the red square, and set smaller than the logo. Use \verb|\epflunit{...}| to typeset one:

\begin{center}
    \epflunit{School of Computer and Communication Sciences}
\end{center}

A branded callout box built on these elements is also available through the \texttt{epflbox} environment:

\begin{epflbox}{Branded callout}
    A white surface with a single Groseille accent and the red square marker, in line with the guideline to balance colors with generous white space.
\end{epflbox}

\subsection{Typography}
\label{sec:brand-typography}
EPFL's official corporate font is \textit{Suisse Int'l} (with Arial as the authorised alternative on office materials). As Suisse Int'l is commercially licensed, this template instead adopts \href{https://vercel.com/font}{Geist}, an open-source (SIL OFL 1.1) sans-serif family in the same Swiss typographic tradition. Three members of the family are used, each where pertinent:

\begin{description}
    \item[Geist Sans] The main document font, used for body text, headings, and the cover pages. {\geistlight Light}, Regular, {\geistmedium Medium}, {\geistsemibold SemiBold}, and \textbf{Bold} weights are available through \verb|\geistlight|, \verb|\geistmedium|, and \verb|\geistsemibold|.
    \item[Geist Mono] Used automatically for all monospaced material: \texttt{file paths}, \verb|inline verbatim|, and code listings rendered with \texttt{minted}.
    \item[Geist Pixel] A decorative display family used for the chapter numerals. It comes in five variants, available as font commands:
\end{description}

\begin{center}
    \setlength{\tabcolsep}{12pt}
    \begin{tabular}{ll}
        {\geistpixelsquare\large EPFL 2026}   & \verb|\geistpixelsquare| \\
        {\geistpixelgrid\large EPFL 2026}     & \verb|\geistpixelgrid| \\
        {\geistpixelcircle\large EPFL 2026}   & \verb|\geistpixelcircle| \\
        {\geistpixeltriangle\large EPFL 2026} & \verb|\geistpixeltriangle| \\
        {\geistpixelline\large EPFL 2026}     & \verb|\geistpixelline| \\
    \end{tabular}
\end{center}

\begin{block}[warning]
Geist Pixel is a display typeface: keep it for headlines, numerals, and accents. Never set body text in it, and never use it where the guidelines require the corporate font.
\end{block}
